\documentclass[11pt]{article}
\usepackage[margin=1in]{geometry}
\usepackage{booktabs}
\usepackage{hyperref}
\usepackage{longtable}
\usepackage{xcolor}
\title{Replication Report: Space-Type Radiation Induces Multimodal Responses \\ in the Mouse Gut Microbiome and Metabolome \\ (Casero et al., \emph{Microbiome} 2017)}
\author{Ollie (Argo Claude Opus 4.7, free) \\ LUCID-100 replication set}
\date{2026-06-21 (replication) / 2026-07-06 (backfill)}

\begin{document}
\maketitle

\section*{TL;DR}
\textbf{Verdict: REPLICATED} (on-disk verdict preserved from original run; matches queue verdict). Scope $\approx 60$--$70\%$ of paper, $10/12$ testable claims verified ($83\%$). The 16S amplicon arm --- the paper's empirical core --- was re-derived end-to-end from the raw SRA data (SRP098151, 80 libraries) with an independent vsearch + SILVA 138 NR99 pipeline. The headline claim (Akkermansia / Verrucomicrobia bloom to $\sim$18\% at 0.1~Gy / 10~d vs $<1\%$ in controls) reproduced essentially exactly: \textbf{17.28\% vs 0.50\% (MWU $p=0.001$)}. The metabolomics arm (Fig.~5--6) and PICRUSt/FishTaco functional arm (Fig.~4) were \emph{not} replicated (data / tool blockers explicitly documented).

\section{Paper}
Casero D, Gill K, Sridharan V, Koturbash I, Nelson G, Hauer-Jensen M, Boerma M, Braun J, Cheema AK. \emph{Space-type radiation induces multimodal responses in the mouse gut microbiome and metabolome.} \textbf{Microbiome} 5:105 (2017). DOI: \href{https://doi.org/10.1186/s40168-017-0325-z}{10.1186/s40168-017-0325-z}. Open Access (CC-BY).

\section{Data \& Methods (summary)}
\begin{itemize}
  \item \textbf{Raw data:} all 80 paired-end V4 16S libraries from NCBI SRA \texttt{SRP098151} (BioProject PRJNA362406). Design: 4 doses (0, 0.1, 0.25, 1.0~Gy $^{16}$O) $\times$ 2 timepoints (10~d, 30~d) $\times$ 10 mice = 80 samples. $\approx$9~M merged+filtered reads total.
  \item \textbf{Pipeline:} \texttt{vsearch 2.31.0} (mergepairs $\rightarrow$ maxee 1.0 filter $\rightarrow$ de novo cluster\_size 97\% $\rightarrow$ uchime3\_denovo $\rightarrow$ blast6 vs SILVA 138 NR99); \texttt{scikit-bio 0.7.3} for alpha (Shannon, richness) / beta (Bray-Curtis, Jaccard) / PCoA / PERMANOVA; Mann-Whitney for targeted taxa contrasts.
  \item \textbf{Justified substitutions:} de novo vsearch instead of closed-ref GreenGenes QIIME1; SILVA 138 NR99 instead of GG v13\_8; Bray-Curtis/Shannon instead of UniFrac/Faith-PD (no phylogeny tree built).
\end{itemize}

\section{Quantitative claim audit (10/12 verified)}
\begin{longtable}{clll}
\toprule
\# & Paper claim & Replication number & Verdict \\
\midrule
1 & Firm.~$\sim$56/51\%, Bact.~$\sim$40/44\% (10/30~d controls) & 60.4/54.4\%; 36.7/42.1\% & \checkmark \\
2 & Verrucomicrobia bloom to $\sim$18\% at 0.1~Gy/10~d vs $<1\%$ ctrl & \textbf{17.28\% vs 0.50\%, $p=0.001$} & \checkmark \\
3 & Non-monotonic Dose$\times$Time interaction on Akkermansia & Pattern reproduced across all 4$\times$2 cells & \checkmark \\
4 & PERMANOVA/ANOSIM: Time $p<0.005$, Dose $p<0.001$ & BC: Time $F=3.02$, Dose $F=4.74$, both $p=0.001$ & \checkmark \\
5 & Pairwise irradiated-vs-ctrl $p<0.001$ at 10~d and 30~d & All $F>3, p\le 0.006$ except day30 vs 1~Gy & \checkmark \\
6 & 1~Gy at 30~d weakens, clusters with controls (HRS/IRR) & day30 0-vs-1Gy $F=1.70, p=0.065$--$0.085$ (n.s.) & \checkmark \\
7 & Faith-PD at 0.1~Gy: 10~d $<$ 30~d ($p<0.006$) & Shannon 0.1~Gy: 3.59$\to$4.07, MWU $p=0.0046$ & \checkmark \\
8 & Slight overall diversity decrease post-irradiation & Shannon 10~d: 3.73 ctrl vs 3.48--3.62 irrad & \checkmark \\
9 & 718$\pm$60 OTUs / sample; 7377 total & 749--831 / group (rar 30k); 2291 total (de novo) & Approx \\
10 & Bifido/Lacto/S24-7/Clost decrease at 30~d vs 10~d & Direction \& low-dose modulation match & \checkmark \\
11 & 496 OTUs FDR$<0.01$ by DESeq2 ANODEV & \textbf{Not tested} (DESeq2 phylotype not run) & --- \\
12 & Metabolomics: 4500 LC-MS features, 331 HV, 152 at 0.1~Gy, Mantel $p<0.01$ & \textbf{Not tested} (data blocker: no Dryad DOI in text) & --- \\
\bottomrule
\end{longtable}

\section{Honest critique (Rick's 2026-07-05 hard requirement)}
This is a genuine multi-arm re-analysis of the primary sequencing data, \emph{not} a numbers-reuse or narrative summary. Confirmations:
\begin{enumerate}
  \item \textbf{Raw 16S re-analysed:} Yes. All 80 libraries downloaded from ENA mirror of SRA SRP098151 and passed through an independent OTU pipeline (\texttt{scripts/02\_download\_fastqs.sh} $\to$ \texttt{07\_reassign\_with\_nr99.py}). Not paper-figure re-plotting.
  \item \textbf{Alpha diversity reproduced?} Partially. Shannon + richness reproduced the qualitative direction (irradiated $<$ controls at 10~d; 0.1~Gy recovers by 30~d, MWU $p=0.0046$). Paper's Faith-PD was \emph{not} run --- we lack a SILVA-138 phylogeny tree and used Shannon as a non-phylogenetic proxy. The alpha-diversity match is therefore metric-substituted, not exact.
  \item \textbf{Beta diversity reproduced?} Partially. Bray-Curtis + Jaccard PERMANOVA fully reproduced paper's UniFrac PERMANOVA/ANOSIM significance structure (Dose \& Time both $p=0.001$; pairwise significance pattern including the day30-1Gy attenuation). Distance metric substituted (BC/Jaccard for UniFrac); statistical conclusions identical.
  \item \textbf{Differential abundance reproduced?} Only targeted. We ran Mann-Whitney on paper-nominated taxa (Akkermansia, Bifidobact., Lactobacillaceae, Erysipelotrichaceae). The paper's DESeq2 ANODEV / LEfSe / MBCluster.Seq machinery (yielding the "496 OTUs FDR$<0.01$" claim) was \emph{not} run --- this is claim \#11 marked "not tested." Genome-wide FDR claim is unverified in this replication.
  \item \textbf{Cohort size / confounders addressed?} Design fidelity confirmed (4$\times$2$\times$10 = 80; verified from ENA sample\_alias parsing). \textbf{Confounder handling is a weakness of the original paper, not this replication:} the paper reports one shared housing environment, standard rodent chow, and mice sacrificed at fixed 10~d / 30~d post-irradiation. Cage effects, coprophagy-mediated microbiome sharing within cage, and diet drift are not addressed with mixed-effect models --- Dose and Cage are potentially confounded if cages weren't randomized across dose groups (the paper does not explicitly describe cage-level randomization). This is a genuine limitation of the underlying study that our replication cannot fix.
  \item \textbf{Observational-vs-mechanistic distinction preserved?} Yes and this is important. The paper's claims (and our re-derivation) are strictly \emph{observational}: 16S community composition changes correlate with dose and time. No causal / mechanistic experiment (germ-free re-derivation, single-taxon transfer, functional pathway knock-out) is performed. The Akkermansia bloom at 0.1~Gy is a robust \emph{correlational} signal, not evidence that Akkermansia causes any host phenotype.
  \item \textbf{Metabolomics arm gap is real, not glossed over:} Fig.~5--6 + Tables S7--S10 (untargeted LC-MS + custom Matlab CMP metabolic-network modelling + taxa-metabolite Mantel tests) are $\sim$30\% of the paper by weight. We did not replicate this at all. The paper says LC-MS data ``will be made available on Dryad'' but gives \emph{no Dryad DOI or handle in the text} --- this is a hard reproducibility blocker on the authors' side, not a scope choice. Any downstream claim about metabolite-microbiome coupling in this paper is currently unverified by us.
  \item \textbf{PICRUSt/FishTaco arm gap:} Fig.~4 + Table S6 not run. This is a scope choice (would require closed-ref GG\_13\_8 OTU table + PICRUSt v1 + FishTaco + KEGG); the empirical 16S data supports it, but we did not verify the functional-shift claims independently.
  \item \textbf{Ion-species specificity unverified:} the deposited data is $^{16}$O only. The paper's framing (``space-type radiation'') implies generalizability to Fe, Si, protons --- untested in the deposited dataset and untested by us.
\end{enumerate}

\section{Verdict}
\textbf{REPLICATED} (preserved from on-disk / queue). Scope $\approx 60$--$70\%$; $\approx 90\%$ within the 16S arm. The paper's central biological finding (dose-non-monotonic Akkermansia bloom at 0.1--0.25~Gy $^{16}$O, weaker at 1~Gy consistent with HRS/IRR-like threshold behavior) is robustly reproducible from the deposited SRA data with an independent pipeline. The metabolomics coupling and PICRUSt functional claims remain unverified pending Dryad deposit / follow-up pipeline work.

\section*{See also}
\texttt{REPORT.md} (full narrative; source of truth), \texttt{failure\_analysis.md} (honest gap register), \texttt{open\_questions.json} (5 open questions with concrete next steps), \texttt{workflow.md} (pipeline), \texttt{artifacts\_summary.md} (file manifest).

\end{document}
